\documentclass[11pt,a4paper]{article}

% ECC Research Laboratory - student report (pdfLaTeX, Overleaf).
\usepackage[T1]{fontenc}
\usepackage[utf8]{inputenc}
\usepackage{lmodern}
\usepackage[margin=2.1cm]{geometry}
\usepackage{microtype}
\usepackage{parskip}
\usepackage{amsmath,amssymb}
\usepackage{booktabs}
\usepackage{tabularx}
\usepackage{array}
\usepackage{float}
\usepackage{needspace}
\usepackage{fancyvrb}
\usepackage{enumitem}
\usepackage{xcolor}
\usepackage{cite}
\renewcommand{\citepunct}{], [}
\renewcommand{\citedash}{]--[}
\usepackage{hyperref}
\hypersetup{colorlinks=true, linkcolor=blue!50!black, urlcolor=blue!60!black, citecolor=blue!50!black,
  pdftitle={Elliptic Curve Cryptography in Practice}, pdfauthor={Diellart Mulolli}}
\setlist{itemsep=0.2em, topsep=0.3em}
\setlength{\emergencystretch}{2em}
\newcolumntype{L}{>{\raggedright\arraybackslash}X}
\newcolumntype{P}[1]{>{\raggedright\arraybackslash}p{#1}}

\title{\textbf{Elliptic Curve Cryptography in Practice}\\[0.3em]
\large From the ECDLP to TLS, Signal, and the post-quantum transition}
\author{Diellart Mulolli (230756200035)\\[0.2em] \small Code and LaTeX source: \url{https://github.com/Diellart-Mulolli/ecc-lab}}
\date{\today}

\begin{document}
\maketitle

\section*{Introduction}
Elliptic curve cryptography (ECC) is today the default public-key technology on the Internet:
almost every TLS 1.3 handshake, every Signal and WhatsApp message, every SSH login with a
modern key and every Bitcoin transaction relies on an elliptic curve. Its attraction is that the
elliptic-curve discrete logarithm problem (ECDLP) admits no known sub-exponential classical
attack, so 256-bit keys give the same security as 3072-bit RSA keys. This report first works through the
mathematics on a toy curve and compares ECC with RSA, then studies how the same curve is used
for key agreement, signatures and encryption, analyses a real-world failure (the 2024 PuTTY
P-521 key-recovery vulnerability), and finally discusses the threat of quantum computers and
how a long-lived system should choose its mechanisms. Three small Python scripts were used to verify the calculations and to measure performance; they are
published, together with the LaTeX source of this report and a running guide, at \url{https://github.com/Diellart-Mulolli/ecc-lab}. Appendix~A lists the values chosen in each
script and its exact output.

%=====================================================================
\section*{Foundations}

\subsection*{Question 1 --- The ECDLP}
Curve (Stamp, Sec.~4.5.1 \cite{stamp}): $E: y^2 \equiv x^3 + 2x + 1 \pmod 5$, base point $P=(0,1)$, $E=\{(0,1),(0,4),(1,2),(1,3),(3,2),(3,3),\infty\}$, so $\#E = 7$ (prime), and $P$ has order $n=7$.

\textbf{(a)} Computation of $2P$ (doubling, $x_1=0,\ y_1=1,\ a=2$):
\begin{align*}
m &= (3x_1^2 + a)(2y_1)^{-1} \bmod 5 = (3\cdot 0 + 2)\cdot 2^{-1} = 2\cdot 3 = 6 \equiv 1,\\
x_3 &= m^2 - x_1 - x_2 = 1 - 0 - 0 = 1, \qquad
y_3 = m(x_1 - x_3) - y_1 = 1\cdot(0-1) - 1 = -2 \equiv 3 .
\end{align*}
So $2P=(1,3)$. Check: $y^2 = 9 \equiv 4$ and $x^3+2x+1 = 1+2+1 = 4$. \checkmark

Computation of $3P = 2P + P = (1,3)+(0,1)$ (addition):
\begin{align*}
m &= (y_2 - y_1)(x_2 - x_1)^{-1} = (1-3)(0-1)^{-1} = (-2)(-1)^{-1} \equiv 3\cdot 4^{-1} = 3\cdot 4 = 12 \equiv 2,\\
x_3 &= m^2 - x_1 - x_2 = 4 - 1 - 0 = 3, \qquad
y_3 = m(x_1 - x_3) - y_1 = 2(1-3) - 3 = -7 \equiv 3 .
\end{align*}
So $3P=(3,3)$. Check: $y^2 = 9 \equiv 4$ and $27+6+1 = 34 \equiv 4$. \checkmark

\textbf{(b)} The private key was found by \emph{exhaustive search}: compute $P, 2P, 3P, \dots$
by repeatedly adding $P$ and compare each result with $Q$:
$2P=(1,3)$, $3P=(3,3)$, $4P = 3P + P = (3,2) = Q$.
Hence $\boxed{k = 4}$. The search needed 3 group operations (one doubling and two additions).
The script \texttt{q1\_toy\_curve.py} reproduces (a) and (b) exactly (Appendix~A.1).
Consistency check: $4P = -3P$ since $7P=\infty$, and indeed $(3,2) = -(3,3)$.
The full cycle is $P=(0,1),\,2P=(1,3),\,3P=(3,3),\,4P=(3,2),\,5P=(1,2),\,6P=(0,4),\,7P=\infty$.
Brute force needs on average $n/2$ operations; this is trivial for $n=7$ but impossible for $n\approx 2^{256}$.

\needspace{8\baselineskip}
\textbf{(c)}
\begin{enumerate}[label=(\roman*)]
\item With the \emph{double-and-add} algorithm (the elliptic-curve analogue of square-and-multiply),
$d$ is processed bit by bit: one doubling per bit and one addition per 1-bit. For a 256-bit $d$
this is about 255 doublings and on average $\approx 128$ additions, i.e.\ fewer than 400 group operations.
Constant-time implementations use the Montgomery ladder (one doubling and one addition per bit,
$\approx 2\times 255$ operations) and fixed-base precomputed tables or windowing (wNAF) to reduce
the number of additions further. Computing $Q=dG$ therefore takes well under a millisecond.
\item The best known classical attack on a well-chosen curve is \emph{Pollard's rho}
(with van Oorschot--Wiener parallelisation). It needs about $\sqrt{\pi n/4}\approx 0.89\sqrt n$
group operations, i.e.\ $\approx 2^{128}$ for P-256. Generic algorithms such as baby-step
giant-step have the same $O(\sqrt n)$ time but also need $O(\sqrt n)$ memory.
This is fully exponential in the key length: 128-bit security.
\item For factoring and for the discrete logarithm in $\mathbb{Z}_p^*$ there exist
\emph{sub-exponential} index-calculus algorithms, above all the General Number Field Sieve,
with cost $L_N[1/3, 1.923] = \exp\big((1.923+o(1))(\ln N)^{1/3}(\ln\ln N)^{2/3}\big)$.
They exploit extra structure (smooth integers, a factor base of small primes). Elliptic-curve points
have no notion of ``small'' or ``smooth'' elements, so index calculus does not apply to
well-chosen curves and only generic $O(\sqrt n)$ attacks remain. Because security grows
linearly with the ECC key length (half the bits) but only like the cube root of the RSA modulus
size, ECC reaches the same security with keys roughly 12 to 30 times shorter.
\end{enumerate}

\textbf{(d)} \emph{Technical reason:} generating a secure curve is hard and easy to get wrong. One
must count the points (Schoof/SEA), ensure that the group order is a large prime (times a small cofactor)
and exclude weak classes (anomalous curves with $\#E=p$, curves with small embedding degree
vulnerable to the MOV/Frey--R\"uck reduction, singular curves, weak twists). Named curves such as P-256
or Curve25519 \cite{sp800186,rfc7748} have been analysed publicly for many years and can be validated against a single
reference. \emph{Practical reason:} interoperability and efficiency. Both peers just send a curve
identifier (e.g.\ \texttt{x25519}, \texttt{secp256r1}) instead of full parameters, and implementations
can use highly optimised, constant-time, formally verified or hardware-accelerated code for a few fixed
curves. Accepting arbitrary parameters has also led to vulnerabilities (CurveBall, CVE-2020-0601,
where Windows accepted attacker-chosen explicit curve parameters).

\subsection*{Question 2 --- ECC versus RSA}
\begin{table}[H]
\textbf{(a)}\par\smallskip
\centering\small
\begin{tabular}{@{}cccc@{}}
\toprule
Security strength & Symmetric & RSA modulus (bits) & ECC key (bits) \\
\midrule
112 & 3TDEA$^{*}$ & 2048 & 224--255 \\
128 & AES-128 & 3072 & 256--383 \\
192 & AES-192 & 7680 & 384--511 \\
256 & AES-256 & 15360 & 512+ \\
\bottomrule
\end{tabular}
\caption{Comparable security strengths, adapted from \cite{nist80057} (Table 2). $^{*}$3TDEA is
deprecated and disallowed after 2023.}
\end{table}

\textbf{(b)} At 128-bit security: $3072/256 = 12$. At 256-bit security: $15360/512 = 30$.
The gap grows because the cost of breaking ECC is fully exponential ($\approx 2^{\ell/2}$ for an
$\ell$-bit key), so doubling security only needs a doubling of the key size, while the GNFS is
sub-exponential: the RSA modulus has to grow roughly with the \emph{cube} of the security level
(up to logarithmic factors). Doubling the security from 128 to 256 bits multiplies the RSA modulus by $5$
(3072 to 15360 bits) but the ECC key only by $2$ (256 to 512 bits).

\begin{table}[H]
\textbf{(c)} Measurements made with \texttt{q2\_benchmark.py} (Python \texttt{cryptography} 41,
OpenSSL backend, one laptop core; average milliseconds per operation; settings and raw output in Appendix~A.2) and sizes from the standards
\cite{rfc8032,fips1865}:\par\smallskip
\centering\small
\begin{tabularx}{\textwidth}{@{}l>{\centering\arraybackslash}X>{\centering\arraybackslash}X>{\centering\arraybackslash}X>{\centering\arraybackslash}X>{\centering\arraybackslash}X@{}}
\toprule
Scheme & Key generation & Signing & Verification & Public key & Signature \\
\midrule
RSA-3072 (PSS) & 813 ms (slowest) & 1.39 ms & \textbf{0.09 ms} & 384 B modulus (422 B DER) & 384 B \\
ECDSA P-256 & 0.06 ms & 0.05 ms & 0.11 ms & 64 B raw (65 B uncompr., 33 B compr.) & 64 B raw (69--72 B DER) \\
Ed25519 & 0.07 ms & 0.06 ms & 0.16 ms & 32 B & 64 B \\
\bottomrule
\end{tabularx}
\caption{RSA-3072 versus ECDSA P-256 versus Ed25519 (128-bit security).}
\end{table}

Key generation in RSA is four orders of magnitude slower (two random 1536-bit primes must be found),
signing is $\approx 25$ times slower, and the keys and signatures are 6--12 times larger.
\emph{RSA is faster at verification}: verification computes $s^e \bmod N$ with the small public exponent
$e = 65537 = 2^{16}+1$, i.e.\ only 17 modular multiplications, whereas ECDSA/EdDSA verification
requires a double scalar multiplication with two full-size 256-bit scalars (several hundred point operations).

\textbf{(d)} For the battery-powered sensors I recommend \textbf{Ed25519} (or ECDSA P-256 if the
hardware only offers a NIST-certified accelerator). The signer is the constrained party: Ed25519
signing is cheap, needs no random number generator per signature (no Sony-style nonce failure on a cheap
device), is easy to implement in constant time, and produces 64-byte signatures, which saves radio
energy compared with 384-byte RSA signatures. The cloud verifies $50\,000/60 \approx 833$ signatures
per second; at $\approx 0.15$ ms per verification this uses about 13\% of one CPU core, and Ed25519 also
supports batch verification. The answer changes partly in the opposite scenario: if a powerful server
signs and millions of small devices only verify, RSA's very cheap verification (small $e$) becomes an
advantage and the large signature size matters less, so RSA-3072 is a reasonable choice (e.g.\ for
firmware images), although an Ed25519 verifier with a 32-byte key is also feasible on most
microcontrollers and much smaller in code and key storage.

%=====================================================================
\clearpage
\section*{Mechanisms and Protocol Scenarios}

\subsection*{Question 3 --- ECDH, ECDSA, EdDSA, and ECIES}
\begin{table}[H]
\textbf{(a)}\par\smallskip
\centering\small
\begin{tabularx}{\textwidth}{@{}P{1.9cm}LLP{2.0cm}P{2.1cm}P{2.3cm}@{}}
\toprule
Mechanism & Security goal & Inputs / outputs & Fresh randomness? & Standard & Real system \\
\midrule
ECDH / X25519 & Key agreement: a shared secret over a public channel (confidentiality, forward secrecy when ephemeral) & In: own private scalar, peer public point. Out: shared point $\to$ KDF $\to$ session keys & Yes, for each ephemeral key pair & RFC 7748, NIST SP 800-56A & TLS 1.3, Signal, WireGuard, SSH \\
ECDSA & Authenticity, integrity, non-repudiation & Sign: $d$, message $\to (r,s)$. Verify: $Q$, message, $(r,s)\to$ valid/invalid & Yes: secret nonce $k$ per signature (unless RFC 6979) & FIPS 186-5, SEC 1, ANSI X9.62 & X.509/TLS certificates, Bitcoin, Ethereum \\
EdDSA / Ed25519 & Authenticity, integrity, non-repudiation & Sign: secret key, message $\to (R,S)$ (64 B). Verify: public key, message, signature $\to$ valid/invalid & No: nonce is derived deterministically & RFC 8032, FIPS 186-5 & OpenSSH, Signal (XEdDSA), Tor onion services, OpenBSD signify \\
ECIES / HPKE & Public-key encryption: confidentiality and integrity to a recipient's public key & In: recipient public key, plaintext, (info, AAD). Out: encapsulated ephemeral key \texttt{enc} + AEAD ciphertext & Yes: a fresh ephemeral key per message & SEC 1, ISO 18033-2; RFC 9180 (HPKE) & TLS ECH, MLS (RFC 9420), Apple iMessage/PQ3, Cloudflare ODoH \\
\bottomrule
\end{tabularx}
\end{table}

\textbf{(b)} The statement is inaccurate because there is no practical elliptic-curve analogue
of the RSA trapdoor permutation $m \mapsto m^e \bmod N$: an EC public key $Q=dG$ is a point,
and there is no efficient way to map an arbitrary message into a point and back with a trapdoor.
(The textbook ``EC-ElGamal'' exists but needs message encoding into points, doubles the ciphertext
size and is malleable, so it is not used.) ECC itself only provides \emph{key agreement} and
\emph{signatures}. In practice, a message is encrypted with a \emph{hybrid} construction, ECIES
or HPKE \cite{rfc9180} (mode \texttt{base}, e.g.\ \texttt{DHKEM(X25519, HKDF-SHA256)} with AES-128-GCM):
\begin{enumerate}
\item \textbf{KEM (key encapsulation):} the sender generates a fresh ephemeral key pair $(e, E=eG)$
and computes $Z = e\,Q_{\text{recipient}}$ (ECDH). $E$ is sent as \texttt{enc}.
\item \textbf{KDF (HKDF):} $Z$, $E$, $Q$ and context information are turned into a uniformly random
symmetric key and nonce.
\item \textbf{DEM / AEAD (AES-GCM or ChaCha20-Poly1305):} encrypts and authenticates the actual data.
\end{enumerate}
The recipient computes $Z = d\,E$, derives the same key and decrypts. ECC handles only the key; the
symmetric cipher handles the data, which also allows arbitrarily long messages.

\textbf{(c)} Both signatures use the same $k$, so $r = x(kG) \bmod n$ is the same. Then
\[
s_1 = k^{-1}(h_1 + r d), \qquad s_2 = k^{-1}(h_2 + r d) \pmod n .
\]
Subtracting eliminates the unknown $d$: $s_1 - s_2 = k^{-1}(h_1 - h_2)$, hence
\[
\boxed{k = (h_1 - h_2)(s_1 - s_2)^{-1} \bmod n}, \qquad
\boxed{d = (s_1 k - h_1)\, r^{-1} \bmod n}.
\]
Both $h_1, h_2$ are hashes of public messages, so everything on the right-hand side is public. The
reuse is visible directly: two signatures with an \emph{identical $r$ value} under the same key
reveal that the same $k$ was used (with a random 256-bit $k$, a repeated $r$ would otherwise happen with
negligible probability). I verified this with \texttt{q3\_nonce\_reuse.py} on P-256 (run with \texttt{-{}-seed 2026}; all chosen
values and the full output are in Appendix~A.3): both signatures share $r = \texttt{0x45f4\ldots 0b8e}$ and are valid,
also when checked with OpenSSL; the two formulas above return exactly the $k$ and $d$ that were used; and the
recovered key forges a valid signature on a new message. In the Sony PS3 case this allowed fail0verflow to sign arbitrary software \cite{ps3}.

\textbf{(d)} \emph{EdDSA} (RFC 8032) hashes the 32-byte secret key with SHA-512; the second half
of the output is a secret \texttt{prefix}. The per-signature value is
$r = \text{SHA-512}(\texttt{prefix} \,\|\, M) \bmod \ell$ and $R = rB$. It is deterministic,
secret and different for every message. \emph{RFC 6979} keeps the ECDSA equations but replaces the
random $k$ by a value generated with HMAC-DRBG seeded by the private key $x$ and $H(m)$, so the same
message always gives the same signature and different messages give independent-looking nonces
\cite{rfc6979}. These designs remove the class of \textbf{randomness failures}: reused, predictable
or biased nonces caused by a weak or broken RNG (Sony PS3, the 2013 Android \texttt{SecureRandom} Bitcoin
thefts). Still possible are \textbf{implementation attacks}: side-channel attacks (timing, cache,
power, EM, e.g.\ Minerva and TPM-Fail which leak nonce bits through non-constant-time scalar multiplication),
\textbf{fault attacks} (deterministic signing is especially vulnerable: a single glitch while signing
the same message twice leaks the key), and incorrect implementations of the derivation itself (PuTTY,
Question 5). Lattice attacks on partially known nonces also remain possible if any nonce bits leak.

\subsection*{Question 4 --- Alice, Bob, and Mallory}
\textbf{(a)} Alice computes $K_A = aB = a(bG)$ and Bob computes $K_B = bA = b(aG)$. Since scalar
multiplication is repeated addition in an abelian group, $a(bG) = (ab)G = (ba)G = b(aG)$, so
$K_A = K_B = abG$. Mallory sees only $A$ and $B$; computing $abG$ from them is the elliptic-curve
Diffie--Hellman problem. The $x$-coordinate is passed through a KDF (HKDF-Extract/Expand in TLS 1.3
\cite{rfc8446,rfc5869}) because (1) it is \emph{not uniformly random}: it is an element of $\mathbb{F}_p$, only about
half of all $x$-values occur on the curve, and it has algebraic structure, while AES keys must be uniform;
(2) one secret must produce \emph{several independent keys} (client/server keys, IVs, handshake versus
application traffic keys, resumption secrets); (3) HKDF binds the key to its \emph{context} (labels and
the transcript hash), so keys cannot be reused across protocols or sessions; and (4) the length must
match (e.g.\ 128 or 256 bits).

\textbf{(b)} Man-in-the-middle attack: Mallory chooses her own secrets $m_1, m_2$.
\begin{itemize}
\item Alice sends $A$; Mallory intercepts it and sends $M_1 = m_1 G$ to Bob as if it came from Alice.
\item Bob sends $B$; Mallory intercepts it and sends $M_2 = m_2 G$ to Alice as if it came from Bob.
\item Alice computes $K_1 = aM_2 = am_2G$ and Mallory computes the same $m_2A$.
Bob computes $K_2 = bM_1 = bm_1G$ and Mallory computes the same $m_1B$.
\end{itemize}
At the end there are two keys: $K_1$ (Alice--Mallory) and $K_2$ (Mallory--Bob). Alice believes she shares
$K_1$ with Bob, Bob believes he shares $K_2$ with Alice; Mallory decrypts, reads or modifies, and re-encrypts every message.
Moving from a 2048-bit finite-field group to a 256-bit curve does not help, because the attack does not
solve any discrete logarithm: Mallory simply performs two honest key exchanges. The problem is
missing \emph{authentication} of the public values, not the hardness of the underlying problem.

\textbf{(c)} In TLS 1.3 the key exchange (\texttt{key\_share} in ClientHello/ServerHello) is
authenticated by the server's \textbf{CertificateVerify} message. It contains a signature (e.g.\
ECDSA P-256 with SHA-256, RSA-PSS, or Ed25519) over 64 bytes of \texttt{0x20}, the context string
``\texttt{TLS 1.3, server CertificateVerify}'', a zero byte, and the \emph{transcript hash} of all handshake
messages from ClientHello up to and including Certificate, which includes both key shares
\cite{rfc8446}. The signature is verified with the public key in the server's leaf \emph{certificate}.
The client trusts that key if the certificate chain leads to a root CA in its trust store, all
signatures in the chain verify, the certificate is valid (dates, key usage), the hostname matches the
Subject Alternative Name, and it is not revoked (and in browsers carries Certificate Transparency
proofs). The \texttt{Finished} MAC then confirms that both sides derived the same keys. If Mallory
replaces $B$, she cannot produce a valid signature over the modified transcript without the server's private key.
\emph{Signal} has no CAs: each user has a long-term identity key (Curve25519/Ed25519 via XEdDSA) and
publishes a \emph{signed prekey} and one-time prekeys to the server. The X3DH/PQXDH handshake mixes
several DH computations involving the identity keys, so only the owner of the identity key can derive
the session keys \cite{x3dh,pqxdh}. Trust in the identity key is obtained by trust on first use plus
optional out-of-band verification of \emph{safety numbers} (comparing 60 digits or scanning a QR code);
a change of identity key is shown to the user as a warning.

\textbf{(d)} \textbf{No.} TLS 1.3 only uses ephemeral (EC)DHE: the session keys came from $abG$, and
$a, b$ were deleted after the handshake. The stolen ECDSA key was used only to \emph{sign}, so it does not
help to compute $abG$; the recorded 2026 sessions remain confidential (\emph{forward secrecy}). With
the stolen key the attacker can \emph{impersonate} the server in the future: run active
man-in-the-middle attacks on new connections or set up a phishing server with the genuine certificate,
until the certificate is revoked and the key replaced. (Exceptions: if session-ticket encryption keys
were also stolen, resumed 0-RTT/PSK-only sessions may be exposed.) In TLS 1.2 with RSA key transport, the client
encrypted the pre-master secret with the server's long-term RSA public key, so the stolen RSA key would
decrypt every recorded pre-master secret and thus \emph{all} past sessions; there is no forward secrecy.
This is one reason TLS 1.3 removed RSA key transport.

%=====================================================================
\section*{Real-World Investigation}

\subsection*{Question 5 --- PuTTY NIST P-521 private-key recovery (CVE-2024-31497)}
\begin{table}[H]
\textbf{(a) Incident profile}\par\smallskip
\centering\small
\begin{tabularx}{\textwidth}{@{}>{\raggedright\arraybackslash}p{3.6cm}L@{}}
\toprule
Discovery and disclosure & Reported privately to the PuTTY developers before disclosure; publicly disclosed on 15 April 2024 together with the fixed release \cite{puttyadv,oss}. \\
Discovered by & Fabian B\"aumer and Marcus Brinkmann, Ruhr University Bochum. \\
Affected products and scale & PuTTY 0.68--0.80 (2017--2024) and Pageant, and products embedding PuTTY code: FileZilla 3.24.1--3.66.5, WinSCP 5.9.5--6.3.2, TortoiseGit 2.4.0.2--2.15.0, TortoiseSVN 1.10.0--1.14.6. Only \emph{NIST P-521} keys (\texttt{ecdsa-sha2-nistp521}) are affected, but every such key used with these tools over 7 years must be considered compromised. PuTTY is one of the most widely used SSH clients on Windows. \\
ECC mechanism & ECDSA signature generation, specifically the derivation of the per-signature nonce $k$ (random number generation). \\
Identifier & CVE-2024-31497; PuTTY advisory \texttt{vuln-p521-bias} \cite{nvd,puttyadv}. \\
Fix & PuTTY 0.81 (15 April 2024) switched all DSA/ECDSA nonce generation to RFC 6979; FileZilla 3.67.0, WinSCP 6.3.3 and TortoiseGit 2.15.0.1 shipped fixes. TortoiseSVN had no patch at disclosure, so users were advised to use Plink from PuTTY 0.81 for SVN over SSH. The fix cannot repair keys that already leaked; they must be revoked and regenerated. \\
\bottomrule
\end{tabularx}
\end{table}

\textbf{(b) Technical root cause.}
\begin{enumerate}
\item Since PuTTY first implemented DSA (around 2001, more than ten years before RFC 6979) it has generated $k$
deterministically, because its author did not trust the entropy source of Windows: $k$ was a SHA-512 hash of
secret data derived from the private key together with a hash of the message, reduced modulo $n$.
\item This is fine for curves with $n < 2^{512}$, but for P-521 the order $n$ has 521 bits, while SHA-512
produces only 512 bits. Reducing a 512-bit value modulo a 521-bit $n$ changes nothing, so the
\textbf{top 9 bits of every nonce are always zero}. The nonce is biased: it is not uniform in $[1,n-1]$.
\item Each signature then gives a linear relation $s_i k_i \equiv h_i + r_i d \pmod n$ with a known small bound
$k_i < 2^{512}$. This is an instance of the \emph{Hidden Number Problem}: with about 60 signatures,
a lattice built from these equations and reduced with LLL/BKZ reveals $d$.
\item The attacker needs about 60 signatures and the public key. They can be collected by (i) a malicious or
compromised SSH server to which the victim logs in (each login produces a signature on data the server knows), or
(ii) passively from public data, e.g.\ Git commits signed through PuTTY/Pageant and published on GitHub.
\end{enumerate}
The failed property is the \textbf{unpredictability} (uniformity) of the nonce, which leads to a failure of
\textbf{authenticity}: the attacker obtains the \textbf{private key}, can log in to every server that trusts it
(supply-chain attacks via Git/SSH access) and can forge signatures in the victim's name.

\textbf{(c) Classification.} The failure is an \textbf{implementation} error: a nonce-generation routine
designed for smaller curves was reused for P-521 without checking that the hash length is at least the bit length of $n$.
The ECDSA standard (FIPS 186-5) and RFC 6979 specify how to generate unbiased nonces; P-521 and the
mathematics of elliptic curves are unaffected. The \textbf{ECDLP was not broken}: the attack never
computes a discrete logarithm. It uses the extra information leaked by the biased nonces together with lattice
reduction. Keys on the same curve generated by correct implementations (e.g.\ OpenSSH) remain secure.

\textbf{(d) Response and lessons.} \emph{Vendors} coordinated disclosure: PuTTY, FileZilla, WinSCP and
TortoiseGit released fixes on the disclosure date or shortly after (TortoiseSVN users were pointed to the patched
Plink in the meantime), and the PuTTY advisory clearly stated that
fixing the software is not enough. Linux distributions shipped patched packages. \emph{Users} were told to
remove the old P-521 public keys from every \texttt{authorized\_keys} file and from GitHub/GitLab, generate new keys
(preferably Ed25519), and re-sign if needed. \emph{Standards}: no standard needed to be withdrawn; the incident
reinforced RFC 6979 and the deterministic design of EdDSA. I consider the technical response adequate and
fast (fixed on the day of disclosure, clear advisory), but the \emph{key rotation} relies on each user, and
signed Git commits remain public forever, so already-leaked keys cannot be ``un-leaked''. \textbf{Proposed control:}
known-answer and statistical tests for nonce generation in the test suite, e.g.\ Project Wycheproof-style tests
that sign many messages for each supported curve and check that the nonces use the full range of $n$ (a
check of the highest bits would have revealed the 9 zero bits immediately). As a process control,
review rule: any reuse of a cryptographic routine for a new parameter set must check its length
assumptions ($|H| \ge \lceil \log_2 n \rceil + 64$).

%=====================================================================
\section*{The Quantum Horizon and Design Decisions}

\subsection*{Question 6 --- ECC, Shor's algorithm, and choosing mechanisms}
\textbf{(a)} Shor's algorithm turns the ECDLP into a \emph{period-finding} problem. For $Q=kP$ in a group of
order $n$, define $f(x,y) = xP + yQ = (x+ky)P$. Then $f(x,y) = f(x',y')$ exactly when
$(x-x') \equiv -k(y-y') \pmod n$, so $f$ is periodic with the hidden period lattice generated by $(k,-1)$ and $(n,0)$.
A quantum computer evaluates $f$ on a superposition of all $(x,y)$ (the costly part: reversible point addition and
modular inversion), then applies the quantum Fourier transform; measurement yields pairs $(u,v)$ with
$u \equiv kv \pmod n$, from which $k = u\,v^{-1} \bmod n$. The cost is polynomial, about $O(\ell^3)$ gates for an $\ell$-bit curve.
\emph{Grover's} algorithm only gives a quadratic speed-up for unstructured search: brute-forcing a 256-bit
key would take $\approx 2^{128}$ quantum steps, which is no better than Pollard rho classically, and Grover
does not parallelise well. So Grover is a concern for symmetric keys (and is handled by AES-256), while
Shor completely breaks ECC.

\emph{Resource estimates.} Litinski (2023) estimates that a 256-bit elliptic-curve private key can be
computed with only about \textbf{50 million Toffoli gates}; on his assumed architecture (6\,000 modules of
1\,152 physical qubits each) one key would be broken about every 10 minutes \cite{litinski}. Earlier work by
Roetteler et al.\ estimated about 2\,330 logical qubits for P-256 \cite{roetteler}. For RSA-2048, Gidney (2025)
estimates that factoring is possible with \textbf{fewer than one million noisy qubits in less than a week},
using about \textbf{6.5 billion Toffoli gates} and roughly 1\,400--1\,550 logical qubits \cite{gidney2025}.
ECC may fall first because the quantum cost depends on the size of the numbers, not on the classical security:
Shor on a 256-bit curve operates on 256-bit numbers, while RSA-2048 needs arithmetic on 2048-bit numbers. The ECC
circuit therefore needs a similar number of logical qubits but about \emph{100 times fewer} gates
($5\cdot 10^{7}$ versus $6.5\cdot 10^{9}$) and much less time. The very property that
makes ECC efficient classically (short keys) makes it a smaller target for a quantum computer.

\textbf{(b)} In a ``harvest now, decrypt later'' (HNDL) attack an adversary records encrypted traffic today
(including the ECDHE key shares) and stores it until a cryptographically relevant quantum computer exists; Shor then
recovers the ephemeral private keys, the session keys and the plaintext. This is urgent for \textbf{ECDH} because
the damage is retroactive: data that must stay confidential for 10--20 years (health, government, trade secrets) is
already at risk, and forward secrecy does not help against Shor. For \textbf{signatures}, the attacker needs the
quantum computer \emph{at the time of verification} to forge: a signature that was verified in a 2026 handshake
cannot be forged retroactively, so migration can wait until shortly before quantum computers arrive.
\emph{Exceptions} that need early migration: signatures that must remain verifiable for a long time or roots of
trust that cannot be changed later, e.g.\ \textbf{firmware/secure-boot keys burned into devices} with a 10--20-year
lifetime, long-lived root CA certificates, and legally binding documents or code signatures that must remain
trustworthy for decades (an attacker could later forge a back-dated document). NIST IR 8547 (draft) plans to deprecate
quantum-vulnerable algorithms at 112-bit security by 2030 and disallow all of them by 2035 \cite{ir8547}.

\begin{table}[H]
\textbf{(c) Design scenario: MedLink (launch 2027, 20-year retention).}\par\smallskip
\centering\small
\begin{tabularx}{\textwidth}{@{}>{\raggedright\arraybackslash}p{2.6cm}>{\raggedright\arraybackslash}p{4.0cm}>{\raggedright\arraybackslash}p{1.7cm}L@{}}
\toprule
Requirement & Mechanism and parameters & PQ / hybrid now? & Justification \\
\midrule
R1. Secure channel app $\leftrightarrow$ server & TLS 1.3 with hybrid key exchange \texttt{X25519MLKEM768} (X25519 + ML-KEM-768, FIPS 203 \cite{fips203}); AES-256-GCM; server authentication with ECDSA P-256 certificate; certificate pinning in the app & \textbf{Yes} (KEM); \textbf{plan} (signature) & Patient data must stay confidential for 20 years, so HNDL applies: the key exchange must be post-quantum now. The hybrid keeps X25519 security if ML-KEM is broken. Authentication is checked online, so the certificate can migrate to ML-DSA later; the app is updatable. \\
R2. Prescriptions verifiable for 10 years & Composite/dual signature ECDSA P-256 (or Ed25519) + ML-DSA-65 (FIPS 204 \cite{fips204}), with RFC 3161 time-stamps and long-term validation & \textbf{Yes} & Signatures must be valid until $\ge 2037$, after the planned 2035 disallowance of ECC; a forged prescription later would be a legal and safety problem. Hybrid gives legal acceptance today (ECDSA) and quantum resistance (ML-DSA). Alternative: SLH-DSA (FIPS 205 \cite{fips205}) for maximum conservatism. \\
R3. Encrypted record export to offline recipient & HPKE (RFC 9180) with a hybrid KEM X25519 + ML-KEM-768 (e.g.\ X-Wing), HKDF-SHA256, AES-256-GCM; record signed by MedLink & \textbf{Yes} & Offline recipient $\Rightarrow$ no interactive handshake, so a one-shot KEM + AEAD is needed. The file can be stored and copied for decades, so it is the classic HNDL target. AES-256 also covers Grover. \\
R4. Firmware updates for a 10-year blood-pressure monitor & Firmware signed with a stateful hash-based signature LMS/HSS (NIST SP 800-208 \cite{sp800208}), optionally in parallel with Ed25519; verification key in boot ROM; anti-rollback counter & \textbf{Yes} & The verification key is fixed in hardware for 10 years and cannot be changed after deployment. LMS verification is only hashing (cheap for a low-power MCU) and its security relies only on SHA-256. State management is easy because MedLink signs few firmware images in an HSM. \\
\bottomrule
\end{tabularx}
\end{table}

\section*{Conclusion}
ECC is the backbone of modern secure communication: the ECDLP has no known sub-exponential classical attack, so
256-bit keys provide 128-bit security with small keys, fast signing and cheap key exchange, which explains its use in
TLS 1.3, Signal, SSH and cryptocurrencies. The calculations on the toy curve showed why computing $kP$ is easy and
recovering $k$ is hard only because $n$ is huge. The case studies show, however, that real-world failures almost never
come from the mathematics: reused or biased nonces (Sony, PuTTY), weak random number generators, side channels and
faulty certificate validation broke real systems while the ECDLP remained intact. Correct, constant-time, deterministic
implementations and named, vetted curves are therefore as important as the choice of the curve itself. Furthermore, ECC
alone does not provide authentication of key exchanges; certificates, signatures or identity verification are required to stop
man-in-the-middle attacks. Its biggest limitation is the quantum threat: Shor's algorithm breaks ECC in polynomial time,
probably with fewer resources than RSA, and recorded ECDH traffic is already exposed to harvest-now-decrypt-later.
The right path for the coming decade is hybrid cryptography (X25519 + ML-KEM, ECDSA + ML-DSA) and crypto-agility, so
that ECC is kept as a well-understood safety net while post-quantum algorithms mature.

%=====================================================================
\clearpage
\section*{Appendix A --- Python Verification: Chosen Values and Outputs}
The three scripts are in the repository \url{https://github.com/Diellart-Mulolli/ecc-lab} (see its README for the
running guide). They were run with Python~3.12.3 and the \texttt{cryptography} package~41.0.7 (OpenSSL backend).
The outputs below are copied unchanged from the terminal. In addition, every value was cross-checked independently:
the toy-curve results by a second computation along a different path, and the P-256 values with OpenSSL
(through \texttt{cryptography}) and the system \texttt{sha256sum} tool.

\subsection*{A.1 \texttt{q1\_toy\_curve.py} (Question 1)}
\begin{table}[H]
\centering\small
\begin{tabularx}{\textwidth}{@{}P{4.2cm}L@{}}
\toprule
Chosen value & Setting in the code \\
\midrule
Curve & $E: y^2 \equiv x^3 + ax + b \pmod p$ with $p = 5$, $a = 2$, $b = 1$ (lab handout, Stamp Sec.~4.5.1) \\
Base point and public key & $P = (0,1)$, $Q = (3,2)$ \\
Point at infinity & represented by \texttt{None}, printed as \texttt{inf} \\
Group law & slope formulas of Stamp, Table~4.1; inverses with \texttt{pow(x, -1, p)} \\
Method for (b) & exhaustive search: add $P$ repeatedly until the result equals $Q$ \\
\bottomrule
\end{tabularx}
\end{table}
Command: \texttt{python3 q1\_toy\_curve.py}
\begin{Verbatim}[fontsize=\footnotesize, frame=single]
Affine points: [(0, 1), (0, 4), (1, 2), (1, 3), (3, 2), (3, 3)] + infinity; #E = 7

(a)
  (0, 1) + (0, 1) (doubling): m = 2 * 2^-1 = 1, x3 = 1^2 - 0 - 0 = 1, y3 = 1(0 - 1) - 1 = 3
  2P = (1, 3) on curve: True
  (1, 3) + (0, 1) (addition): m = 3 * 4^-1 = 2, x3 = 2^2 - 1 - 0 = 3, y3 = 2(1 - 3) - 3 = 3
  3P = (3, 3) on curve: True

(b) brute force for Q = (3, 2)
  2P = (1, 3)
  3P = (3, 3)
  4P = (3, 2)
  k = 4, found after 3 point additions/doublings

All multiples of P:
  1P = (0, 1)
  2P = (1, 3)
  3P = (3, 3)
  4P = (3, 2)
  5P = (1, 2)
  6P = (0, 4)
  7P = inf
  8P = (0, 1)
\end{Verbatim}
\textbf{Check.} The six affine points match the set in the lab handout, and $\#E = 7$. The slopes $m = 1$ and $m = 2$,
the points $2P = (1,3)$ and $3P = (3,3)$, and $k = 4$ after 3 operations are the hand results of Question~1(a)--(b).
$7P = \infty$ confirms that $P$ has order 7. Independent cross-check: computing $4P$ as the doubling $2(2P)$ instead of
$3P + P$ also gives $(3,2)$, and counting points with Euler's criterion also gives $\#E = 7$.

\subsection*{A.2 \texttt{q2\_benchmark.py} (Question 2(c))}
\begin{table}[H]
\centering\small
\begin{tabularx}{\textwidth}{@{}P{4.2cm}L@{}}
\toprule
Chosen value & Setting in the code \\
\midrule
Signed message & \texttt{"sensor reading:\ 21.4 C, 48 \% RH"} (31 bytes) \\
RSA & 3072-bit modulus, $e = 65537$, RSA-PSS with MGF1-SHA-256, 32-byte salt, SHA-256 \\
ECDSA & curve P-256 (\texttt{secp256r1}), SHA-256, DER-encoded signature \\
EdDSA & Ed25519 (RFC 8032, SHA-512 internally) \\
Repetitions & key generation: 5 for RSA, 500 for ECDSA and Ed25519; signing and verification: 300 each \\
Measurement & \texttt{time.perf\_counter()}, mean time per operation in ms; key sizes as X.509 DER (\texttt{SubjectPublicKeyInfo}) \\
\bottomrule
\end{tabularx}
\end{table}
Command: \texttt{python3 q2\_benchmark.py} (the run used for Table~2)
\begin{Verbatim}[fontsize=\footnotesize, frame=single]
scheme            keygen ms   sign ms  verify ms  pk DER B  sig B
RSA-3072 (PSS)      813.091     1.386      0.092       422    384
ECDSA P-256           0.058     0.047      0.114        91     72
Ed25519               0.067     0.060      0.155        44     64

P-256 public key: uncompressed 65 B, compressed 33 B
Ed25519 public key raw: 32 B
\end{Verbatim}
\textbf{Check.} Table~2 rounds these numbers. Timings change from run to run (three runs gave RSA key generation
times of 356, 813 and 1086~ms), but the order is always the same: RSA is slowest at key generation and signing and
fastest at verification. Derived values: $813.091/0.058 \approx 14\,000$ (four orders of magnitude),
$1.386/0.047 \approx 29$ and $1.386/0.060 \approx 23$ (signing $\approx 25\times$ slower), and
$50\,000/60 \cdot 0.155~\text{ms} \approx 0.13$~s of CPU time per second ($\approx 13\%$). The sizes were confirmed
separately: RSA-3072 modulus and signature 384~B, DER public keys 422~B (RSA), 91~B (P-256) and 44~B (Ed25519).
DER-encoded ECDSA signatures varied between 69 and 72~B over 200 signatures, because DER drops leading zero bytes of
$r$ and $s$ and adds a sign byte when the top bit is set.

\clearpage
\subsection*{A.3 \texttt{q3\_nonce\_reuse.py} (Question 3(c))}
\begin{table}[H]
\centering\small
\begin{tabularx}{\textwidth}{@{}P{3.4cm}L@{}}
\toprule
Chosen value & Setting in the code \\
\midrule
Curve & NIST P-256 \cite{fips1865,sp800186}: $p = 2^{256} - 2^{224} + 2^{192} + 2^{96} - 1$, $a = -3$ \\
\quad $p$ & \texttt{0xffffffff00000001000000000000000000000000ffffffffffffffffffffffff} \\
\quad order $n$ & \texttt{0xffffffff00000000ffffffffffffffffbce6faada7179e84f3b9cac2fc632551} \\
\quad $G_x$ & \texttt{0x6b17d1f2e12c4247f8bce6e563a440f277037d812deb33a0f4a13945d898c296} \\
\quad $G_y$ & \texttt{0x4fe342e2fe1a7f9b8ee7eb4a7c0f9e162bce33576b315ececbb6406837bf51f5} \\
Hash & $h = \text{SHA-256}(m)$ read as an integer, reduced mod $n$ \\
Messages & $m_1 = \texttt{"firmware 3.55"}$, $m_2 = \texttt{"homebrew loader"}$ \\
Private key $d$, nonce $k$ & uniform in $[1, n-1]$ from Python's \texttt{random.Random(2026)} (\texttt{-{}-seed 2026}, for reproducibility only; without the option the script uses the cryptographic \texttt{secrets} module) \\
Scalar multiplication & left-to-right double-and-add (demonstration only, not constant-time) \\
Forgery & message \texttt{"anything the attacker wants"} signed with the recovered $d$ and nonce 12345 \\
\bottomrule
\end{tabularx}
\end{table}
Command: \texttt{python3 q3\_nonce\_reuse.py -{}-seed 2026}
\begin{Verbatim}[fontsize=\footnotesize, frame=single]
victim private key d = 0xf3f49249dc28ff90a5aec7978306d03bf38b2ffc80a4df5a51c9bc701e7ea41a
reused nonce       k = 0x6bad6be28e7aa6e99f19950499dd251de512148239292d22e255accb1a466885
signature 1 on 'firmware 3.55':
  h1 = 0xba2c60557cd36927df60802fa63d01ff7bccb978daee9847b797aedbafcd7bbc
  r1 = 0x45f4600c8a058ddaea76061d9eb761b67b7799eb14fbf231d21bf31f95c60b8e
  s1 = 0xa22af549b244c43a15a430634a4d87ff55eb0dcbc0ba4c7bd8e1eaa53b1ea7b8
signature 2 on 'homebrew loader':
  h2 = 0xaabb15092f5617c4891d2a5517315f2069371074a2af28a33484e2c0de2576b2
  r2 = 0x45f4600c8a058ddaea76061d9eb761b67b7799eb14fbf231d21bf31f95c60b8e
  s2 = 0x33ccd118c13f82ff4ed22302c5f1b2477997f28853e00b7295e2aca1da1edb11
both signatures valid: True True
r1 == r2 (visible reuse): True
recovered k = 0x6bad6be28e7aa6e99f19950499dd251de512148239292d22e255accb1a466885
recovered d = 0xf3f49249dc28ff90a5aec7978306d03bf38b2ffc80a4df5a51c9bc701e7ea41a
recovered k correct: True
recovered d correct: True
forged signature verifies: True
\end{Verbatim}
\textbf{Check.} $r_1 = r_2$ exposes the reused nonce. Inserting the printed values into the formulas of Question~3(c):
\begin{align*}
k &= (h_1 - h_2)(s_1 - s_2)^{-1} \bmod n\\
  &= \text{\small\texttt{0x6bad6be28e7aa6e99f19950499dd251de512148239292d22e255accb1a466885}},\\
d &= (s_1 k - h_1)\, r^{-1} \bmod n\\
  &= \text{\small\texttt{0xf3f49249dc28ff90a5aec7978306d03bf38b2ffc80a4df5a51c9bc701e7ea41a}},
\end{align*}
which are exactly the nonce and private key chosen by the script. Independent cross-checks with OpenSSL: the
generator $G$ in the script equals OpenSSL's P-256 generator; $x(kG) \bmod n$ computed by OpenSSL equals $r$; both
signatures are accepted by OpenSSL's ECDSA verification under $Q = dG$; and $h_1$, $h_2$ equal the
\texttt{sha256sum} digests of the two messages.

\clearpage
\begin{thebibliography}{99}
\footnotesize\setlength{\itemsep}{1pt}\setlength{\parsep}{0pt}\setlength{\parskip}{0pt}
\bibitem{stamp}
M.~Stamp, \emph{Information Security: Principles and Practice}, 3rd~ed. Hoboken, NJ, USA: Wiley, 2022.

\bibitem{sp800186}
L.~Chen, D.~Moody, K.~Randall, A.~Regenscheid, and A.~Robinson, ``Recommendations for discrete logarithm-based cryptography: Elliptic curve domain parameters,'' Nat.\ Inst.\ Standards Technol., Gaithersburg, MD, USA, NIST SP~800-186, Feb.~2023, doi: \href{https://doi.org/10.6028/NIST.SP.800-186}{10.6028/NIST.SP.800-186}.

\bibitem{rfc7748}
A.~Langley, M.~Hamburg, and S.~Turner, ``Elliptic curves for security,'' Internet Eng.\ Task Force, RFC~7748, Jan.~2016. [Online]. Available: \url{https://www.rfc-editor.org/rfc/rfc7748}

\bibitem{nist80057}
E.~Barker, ``Recommendation for key management: Part~1 -- General,'' Nat.\ Inst.\ Standards Technol., Gaithersburg, MD, USA, NIST SP~800-57 Part~1 Rev.~5, May 2020, doi: \href{https://doi.org/10.6028/NIST.SP.800-57pt1r5}{10.6028/NIST.SP.800-57pt1r5}.

\bibitem{rfc8032}
S.~Josefsson and I.~Liusvaara, ``Edwards-curve digital signature algorithm (EdDSA),'' Internet Eng.\ Task Force, RFC~8032, Jan.~2017. [Online]. Available: \url{https://www.rfc-editor.org/rfc/rfc8032}

\bibitem{fips1865}
\emph{Digital Signature Standard (DSS)}, FIPS PUB~186-5, Nat.\ Inst.\ Standards Technol., Gaithersburg, MD, USA, Feb.~2023, doi: \href{https://doi.org/10.6028/NIST.FIPS.186-5}{10.6028/NIST.FIPS.186-5}.

\bibitem{rfc9180}
R.~Barnes, K.~Bhargavan, B.~Lipp, and C.~Wood, ``Hybrid public key encryption,'' Internet Eng.\ Task Force, RFC~9180, Feb.~2022. [Online]. Available: \url{https://www.rfc-editor.org/rfc/rfc9180}

\bibitem{ps3}
fail0verflow, ``Console hacking 2010: PS3 epic fail,'' presented at the 27th Chaos Commun.\ Congr.\ (27C3), Berlin, Germany, Dec.~2010.

\bibitem{rfc6979}
T.~Pornin, ``Deterministic usage of the digital signature algorithm (DSA) and elliptic curve digital signature algorithm (ECDSA),'' Internet Eng.\ Task Force, RFC~6979, Aug.~2013. [Online]. Available: \url{https://www.rfc-editor.org/rfc/rfc6979}

\bibitem{rfc8446}
E.~Rescorla, ``The transport layer security (TLS) protocol version~1.3,'' Internet Eng.\ Task Force, RFC~8446, Aug.~2018. [Online]. Available: \url{https://www.rfc-editor.org/rfc/rfc8446}

\bibitem{rfc5869}
H.~Krawczyk and P.~Eronen, ``HMAC-based extract-and-expand key derivation function (HKDF),'' Internet Eng.\ Task Force, RFC~5869, May 2010. [Online]. Available: \url{https://www.rfc-editor.org/rfc/rfc5869}

\bibitem{x3dh}
M.~Marlinspike and T.~Perrin, ``The X3DH key agreement protocol,'' Signal Messenger, Nov.~2016. Accessed: Oct.~13, 2026. [Online]. Available: \url{https://signal.org/docs/specifications/x3dh/}

\bibitem{pqxdh}
E.~Kret and R.~Schmidt, ``The PQXDH key agreement protocol,'' Signal Messenger, 2023. Accessed: Oct.~13, 2026. [Online]. Available: \url{https://signal.org/docs/specifications/pqxdh/}

\bibitem{puttyadv}
S.~Tatham, ``PuTTY vulnerability vuln-p521-bias,'' PuTTY wishlist, Apr.~2024. Accessed: Oct.~13, 2026. [Online]. Available: \url{https://www.chiark.greenend.org.uk/~sgtatham/putty/wishlist/vuln-p521-bias.html}

\bibitem{oss}
F.~B\"aumer and M.~Brinkmann, ``CVE-2024-31497: Secret key recovery of NIST P-521 private keys through biased ECDSA nonces in PuTTY client,'' oss-security mailing list, Apr.~15, 2024. Accessed: Oct.~13, 2026. [Online]. Available: \url{https://www.openwall.com/lists/oss-security/2024/04/15/6}

\bibitem{nvd}
Nat.\ Inst.\ Standards Technol., ``CVE-2024-31497 detail,'' National Vulnerability Database. Accessed: Oct.~13, 2026. [Online]. Available: \url{https://nvd.nist.gov/vuln/detail/CVE-2024-31497}

\bibitem{litinski}
D.~Litinski, ``How to compute a 256-bit elliptic curve private key with only 50 million Toffoli gates,'' 2023, \emph{arXiv:2306.08585}.

\bibitem{roetteler}
M.~Roetteler, M.~Naehrig, K.~M. Svore, and K.~Lauter, ``Quantum resource estimates for computing elliptic curve discrete logarithms,'' in \emph{Advances in Cryptology -- ASIACRYPT 2017} (Lecture Notes in Computer Science, vol.~10625). Cham, Switzerland: Springer, 2017, pp.~241--270, doi: \href{https://doi.org/10.1007/978-3-319-70697-9_9}{10.1007/978-3-319-70697-9\_9}.

\bibitem{gidney2025}
C.~Gidney, ``How to factor 2048 bit RSA integers with less than a million noisy qubits,'' 2025, \emph{arXiv:2505.15917}.

\bibitem{ir8547}
D.~Moody, R.~Perlner, A.~Regenscheid, A.~Robinson, and D.~Cooper, ``Transition to post-quantum cryptography standards,'' Nat.\ Inst.\ Standards Technol., Gaithersburg, MD, USA, NIST IR~8547 (Initial Public Draft), Nov.~2024, doi: \href{https://doi.org/10.6028/NIST.IR.8547.ipd}{10.6028/NIST.IR.8547.ipd}.

\bibitem{fips203}
\emph{Module-Lattice-Based Key-Encapsulation Mechanism Standard}, FIPS PUB~203, Nat.\ Inst.\ Standards Technol., Gaithersburg, MD, USA, Aug.~2024, doi: \href{https://doi.org/10.6028/NIST.FIPS.203}{10.6028/NIST.FIPS.203}.

\bibitem{fips204}
\emph{Module-Lattice-Based Digital Signature Standard}, FIPS PUB~204, Nat.\ Inst.\ Standards Technol., Gaithersburg, MD, USA, Aug.~2024, doi: \href{https://doi.org/10.6028/NIST.FIPS.204}{10.6028/NIST.FIPS.204}.

\bibitem{fips205}
\emph{Stateless Hash-Based Digital Signature Standard}, FIPS PUB~205, Nat.\ Inst.\ Standards Technol., Gaithersburg, MD, USA, Aug.~2024, doi: \href{https://doi.org/10.6028/NIST.FIPS.205}{10.6028/NIST.FIPS.205}.

\bibitem{sp800208}
D.~A. Cooper, D.~C. Apon, Q.~H. Dang, M.~S. Davidson, M.~J. Dworkin, and C.~A. Miller, ``Recommendation for stateful hash-based signature schemes,'' Nat.\ Inst.\ Standards Technol., Gaithersburg, MD, USA, NIST SP~800-208, Oct.~2020, doi: \href{https://doi.org/10.6028/NIST.SP.800-208}{10.6028/NIST.SP.800-208}.
\end{thebibliography}

\end{document}
